\documentclass[12pt]{article}
\usepackage[utf8]{inputenc}
\usepackage[T1]{fontenc}
\usepackage{hyperref}
\title{Zhuang, Lianshan Language Resources for K-12 Curriculum Development}
\author{Compiled by Claude Sonnet 5.5 (claude-sonnet-5-5) for Andr\'as Kornai}
\date{}
\begin{document}
\maketitle

\section{Introduction}
Lianshan Zhuang (ISO 639-3: zln; Glottocode lian1252) is a Tai-Kadai language of the Zhuang macrolanguage, spoken in Lianshan and northern Huaiji, Guangdong, China. Speaker counts differ by source (about 33,000 to 48,000). It is reported as taught in no schools. Searches found no primers, textbooks, dictionaries or curricula specific to Lianshan Zhuang. This dossier lists only the reference and catalogue entries that were found.

\section{Primary Educational Resources}
No primers, textbooks or school programmes specific to Lianshan Zhuang were found in the search results.

\section{Linguistic and Reference Materials}
\begin{itemize}
\item Glottolog entry for the language (Lianshan-Liujiang group): \url{https://glottolog.org/resource/languoid/id/lian1252.bigmap.html}
\item Overview of the Zhuang languages, including writing systems (Sawndip, standard Zhuang): \url{https://en.wikipedia.org/wiki/Zhuang_languages}
\item ISO 639:z listing: \url{https://en.wikipedia.org/wiki/ISO_639:z}
\end{itemize}

\section{Religious and Community Resources}
\begin{itemize}
\item Joshua Project people-group profile: \url{https://joshuaproject.net/people_groups/19800}
\end{itemize}
No community organisations were found.

\section{Digital and Audio Resources}
\begin{itemize}
\item Global Recordings Network page for zln, which lists no recordings for this entry: \url{https://globalrecordings.net/en/language/zln}
\item Global Recordings Network page for a related ``Liansan'' variety, which reportedly offers Bible-story audio. The match to Lianshan was not verified: \url{https://www.globalrecordings.net/en/language/6136}
\end{itemize}

\section{Curriculum Development Context}
The language is used mainly by older adults, and no school use is known. Any curriculum would have to begin with documentation: recording speakers, building a word list and settling on a writing system. Standard Zhuang materials may serve as a regional reference, but they are not Lianshan materials. The search results did not show how closely the two are related.

\end{document}